\begin{figure}[t]
\centering
\begin{tikzpicture}
\begin{axis}[
  width=0.78\columnwidth,
  height=0.48\columnwidth,
  xlabel={Compactness $C$},
  ylabel={$\mathcal R_{\rm N}=\sqrt{\lambda_{\rm soft}/\lambda_{\rm hard}}$},
  ymode=log,
  ymin=1e-3,ymax=1e-1,
  xmin=0.06,xmax=0.18,
  grid=both,
  legend style={draw=none,at={(0.03,0.97)},anchor=north west},
]
\addplot[thick,mark=*] table[x=C,y=RN]{data/normal_spectrum.dat};
\addlegendentry{normal response}
\end{axis}
\begin{axis}[
  width=0.78\columnwidth,
  height=0.48\columnwidth,
  xmin=0.06,xmax=0.18,
  ymin=0.990,ymax=0.997,
  axis x line=none,
  axis y line*=right,
  ylabel={I--Love alignment},
  legend style={draw=none,at={(0.97,0.03)},anchor=south east},
]
\addplot[thick,dashed,mark=square*] table[x=C,y=alignment]{data/normal_spectrum.dat};
\addlegendentry{$|n_{\rm soft}\!\cdot n_{I{\rm Love}}|$}
\end{axis}
\end{tikzpicture}
\caption{Normal-response hierarchy along the $\Gamma=2$ reference sequence. After quotienting the stellar-sequence tangent, the soft-to-hard normal RMS ratio remains well below unity (left axis), while the soft covector stays closely aligned with the embedded I--Love normal (right axis).}
\label{fig:normal-spectrum}
\end{figure}
